\chapter{Software and reproducibility}
\label{chap:08-logiciel}

In 2018, Pineau asked for work that is reproducible, reusable and robust. This chapter describes how
the project repository, \repo, tries to meet these three demands. It presents the layout of the
repository, the Python package that gathers the reusable computations, the file formats, the
mechanical check of every published number, the tests, continuous integration, the frozen
environments, the hardware and the compute used. A last section states what the repository lacks and
how to obtain it.

\section{Repository layout}
\label{sec:depot}

The repository follows a one-way flow, from runs to documents
(\cref{fig:architecture,fig:architecture-2}); \cref{tab:depot} names its directories. A run writes
a raw result file: an evaluation curve in CSV format for a simulation, a JSONL log for a GRPO run, a
JSON reading file for a paper of the survey. An analysis script reads these files and writes a JSON
result file. Two checking scripts then recompute every number of the texts from these JSON files.
Every number in a document comes from them.

\fig{architecture}{The repository, first part: from runs to JSON files. Training runs write the raw results (CSV, JSONL, logs). The analysis scripts, built on the \file{pineau\_rl} package, derive one JSON file per experiment from them. \texttt{make analysis} redoes this whole chain without any training.}

\fig{architecture-2}{The repository, second part: from JSON files to documents. \file{check\_numbers.py} and \file{check\_article.py} hold every number of the texts against its JSON file, \file{tools/tables.py} generates the tables of this report. \texttt{make verify} chains analyses, checks and tests, then confirms that every JSON file comes back identical. Continuous integration reruns this path at every push to GitHub.}

\begin{table}[tbp]
\centering
\small
\caption{The directories of the repository and their role.}
\label{tab:depot}
\begin{tabularx}{\linewidth}{@{}lX@{}}
\toprule
Path & Content \\
\midrule
\file{study/} & The survey: protocol (\file{PROTOCOL.md}), reader's sheet (\file{GRILLE.md}),
  draws, readings (\file{codes/}, \file{codes\_v2/}, \file{codes\_x8/}), harmonisation
  (\file{overrides*.json}), provenance of the readers (\file{readers/}), analysis scripts and JSON result
  files, among them those of deviations 30 to 36 (\file{survey\_extra.py}, \file{x8\_stats.py},
  \file{q5\_versions\_stats.py}, \file{gemma\_stats.py}, \file{r36\_stats.py}). \\
\file{sim/} & The simulations: protocol (\file{SIM\_PROTOCOL.md}), one run (\file{run.py}), a
  re-evaluation under modified physics (\file{dose.py}), run lists (\file{jobs*.txt}), launchers and
  watchdog (\file{drive.sh}, \file{queue.sh}, \file{worker.sh}, \file{scale.sh},
  \file{watchdog.sh}), curves (\file{results/}, with the two evaluation videos of E6 in
  \file{results/e6\_video/}), logs (\file{logs/}), final policies (\file{models/}), the agent published
  by the Zoo and its files (\file{zoo\_hub/}). \\
\file{sim/analysis/} & One script and one JSON file per experiment (E1 to E7, B1, B3, dose,
  Bonferroni correction), with their variants at twenty random seeds and at equal budget. The analyses
  of deviations 21 to 31, one per file: the strict count of E2 (\file{e2\_strict.py}), the \task{Hopper}
  reversal tested (\file{e7\_test.py}), E1 on sixteen pairs, E3 with Holm, E5 out of sample, the
  calibration of the rule, the factor of B1, the curves, the effect sizes, the speed of the E6 agents. The
  survey of the hyperparameters (\file{hyperparams.py}). Three analyses written in advance, waiting for
  their training runs (\file{c6\_determinism.py}, \file{e6\_gamma.py}, \file{c2.py}). \\
\file{sim/llm/} & The GRPO strand: protocol, \file{grpo.py}, launchers, diagnostics
  (\file{diag/}), logs, results, \file{summary.py} and \file{logs\_stats.py} (deviation~19). The rerun of
  deviation~20: its launcher \file{c7.sh} and its analysis \file{c7\_stats.py}, written before any run. \\
\file{src/pineau\_rl/} & The Python package of reusable computations (\cref{sec:paquet}). \\
\file{tests/} & The unit tests (\cref{sec:tests}). \\
\file{envs/} & The frozen environments, the platform and the document toolchain
  (\cref{chap:D-environnements}). \\
\file{docs/} & The summary document in two languages, its bibliography, its figures and
  \file{check\_numbers.py}. \\
\file{article/} & The blog article in two languages, its figures and \file{check\_article.py}. \\
\file{report/} & This technical report, its generated tables and its figures. \\
\file{tools/} & \file{compare\_results.py}, which compares JSON files value by value. \\
\bottomrule
\end{tabularx}
\end{table}

\section{The \texttt{pineau\_rl} package}
\label{sec:paquet}

The computations that serve several scripts form an installable Python package,
\file{src/pineau\_rl/}, version 1.0.0, under the BSD 3-Clause licence. Its only dependencies are
NumPy, SciPy, statsmodels and rliable. Every analysis script imports it. The \file{Makefile} puts
\file{src/} on Python's search path, so the scripts run without installing the package. Submodules
load on demand: \texttt{import pineau\_rl.stats} does not load rliable. The package has four public
modules (\texttt{stats}, \texttt{bootstrap}, \texttt{runs}, \texttt{grid}) and a formatting module
(\texttt{fmt}).

\subsection{\texttt{pineau\_rl.stats}}

This module gathers the statistics of the survey and of the experiments. Every function is
deterministic. The module header warns that the computations follow the exact order of the original
scripts, so that the result files regenerate byte for byte: simplifying an expression changes the
last digit of a result.
\begin{itemize}
  \item \texttt{wilson(k, n, z=1.96)} and \texttt{wilson\_ci(k, n, alpha=0.05)}: the interval of
        \textcite{wilson1927probable} for a proportion $k/n$. With $\hat p = k/n$, its centre is
        $(\hat p + z^2/2n)/(1 + z^2/n)$ and its half-width
        $z\sqrt{\hat p(1-\hat p)/n + z^2/4n^2}\,/(1 + z^2/n)$. The second function delegates to
        statsmodels, with $z = \Phi^{-1}(1-\alpha/2)$.
  \item \texttt{newcombe\_diff(ka, na, kb, nb)}: the 95\% interval for the difference of two
        proportions by the hybrid method of \textcite{newcombe1998interval}, built from the two
        Wilson intervals $(l_a, u_a)$ and $(l_b, u_b)$:
        $d - \sqrt{(\hat p_a - l_a)^2 + (u_b - \hat p_b)^2}$ and
        $d + \sqrt{(u_a - \hat p_a)^2 + (\hat p_b - l_b)^2}$, with $d = \hat p_a - \hat p_b$.
  \item \texttt{fisher(ka, na, kb, nb)}: the two-sided $p$-value of Fisher's exact test on the
        $2 \times 2$ table.
  \item \texttt{compare(a, b)}: the two previous functions combined, \texttt{\{diff, lo, hi, p\}}, at
        full precision.
  \item \texttt{holm(pvalues)}: the $p$-values adjusted by the step-down procedure of
        \textcite{holm1979simple}, $\tilde p_{(i)} = \min\bigl(1, \max_{j \le i}(m-j+1)\,p_{(j)}\bigr)$.
  \item \texttt{cochran\_armitage(k, n, x)}: the test for a linear trend in proportions
        \parencite{armitage1955tests}. With $\bar p = \sum k_i / \sum n_i$ and
        $\bar x = \sum x_i n_i / \sum n_i$, the statistic is
        $z = \sum x_i (k_i - n_i \bar p) \big/ \sqrt{\bar p(1-\bar p)\sum n_i (x_i - \bar x)^2}$.
  \item \texttt{cohen\_kappa(a, b)}: the kappa of \textcite{cohen1960coefficient},
        $(p_o - p_e)/(1 - p_e)$, \texttt{None} when $p_e = 1$.
  \item \texttt{probability\_of\_improvement(x, y)}: $P(X > Y) + \tfrac12 P(X = Y)$ over all pairs
        of runs, which is the Mann-Whitney $U$ statistic divided by $n_x n_y$
        \parencite{agarwal2021precipice}.
  \item \texttt{welch\_curve(a, b)}, \texttt{welch\_pooled(a, b)} and
        \texttt{henderson\_script(M, S, a, b, n)}: the three readings of the statistic of
        \textcite{henderson2018matters} studied by experiment B1. The first applies the test of
        \textcite{welch1947generalization} at each evaluation point, then averages the $t$ values or
        the $p$-values. The second stacks all the points of the curves as independent observations.
        The third redoes the computation of the script \file{from\_same\_dist.py}
        \parencite{henderson2017samedist}: a Welch test built from the means and standard deviations
        across episodes, into which the spread between random seeds never enters.
\end{itemize}

\begin{lstlisting}[language=Python,caption={The computation of Henderson's script, in \texttt{pineau\_rl.stats}.},captionpos=b]
def henderson_script(M, S, a, b, n):
    t, p = _st.ttest_ind_from_stats(M[a].mean(0), S[a].mean(0), n,
                                    M[b].mean(0), S[b].mean(0), n, equal_var=False)
    return t.mean(), p.mean()
\end{lstlisting}

\subsection{\texttt{pineau\_rl.bootstrap}}

The bootstrap module wraps the stratified bootstrap intervals of rliable \parencite{agarwal2021precipice}.
\texttt{iqm(x)} returns the interquartile mean, the mean of the middle half of the scores.
\texttt{iqm\_interval(scores, reps=2000)} returns, for each group of runs, the interquartile mean and
its 95\% percentile bootstrap interval. \texttt{improvement\_interval(x, y, reps=2000,
confidence=0.95)} returns the probability of improvement and its interval, at the requested level:
the Bonferroni analysis of deviation 18 of the simulations uses it with $1 - 0.05/24$.
\texttt{probability\_of\_improvement\_bootstrap(x, y, rng)} resamples $x$ and $y$ with an explicit
NumPy generator, without rliable.

A defect of rliable 1.2.0 dictates how this module is used. The library accepts a
\texttt{random\_state} argument but draws its resamples from NumPy's global generator (deviation 13
of the simulations). A script therefore seeds this generator once, at the start
(\texttt{np.random.seed(2026)}), then calls the functions in a fixed order: each call consumes the
global generator, so an interval depends on the calls that precede it. The module imports rliable
as soon as the module itself loads, never inside a function, so that no import happens after the
random seed is set.

\subsection{\texttt{pineau\_rl.runs}}

The runs module defines the format of a simulation run (\cref{sec:formats}) and its readers.
\texttt{run\_path(root, algo, task, seed)} gives the path of a run. \texttt{read\_run(path,
final\_step=None, upto=None)} reads a file and checks it: required columns present, steps strictly
increasing. It returns \texttt{None} for a missing file or an unfinished run. The argument
\texttt{upto} cuts the curve at a smaller budget, which serves the equal-budget analysis.
\texttt{write\_run(path, rows)} writes the format from any library. \texttt{finals(\dots)} returns the
score of each random seed at the requested step, \texttt{load\_grid(\dots)} the full grid
\texttt{\{task: \{algorithm: scores\}\}} with the list of missing pairs.

\subsection{\texttt{pineau\_rl.grid}}

The grid module reads the survey grid. It defines the lists of criteria (\texttt{V1}: Q1 to Q8 and X1 to
X7; \texttt{V2}: Y1 to Y12), the integer criteria (\texttt{INT}: X1, X6, Y3), the allowed values and
the kinds of uncertainty of Y11 (\texttt{KINDS}). \texttt{load\_reads(root, readers)} loads the
reading files, \texttt{load\_overrides(path)} the harmonisation. The final value of a criterion
follows one of two semantics, kept apart because the scripts differ: \texttt{value\_v1} (one reading
and its harmonisation) and \texttt{value} (the reading of the v2 grid when it has the criterion,
otherwise that of the v1 grid, with Y3 turned into yes or no from 2 upward). \texttt{tally(values)}
returns the numerator and the denominator of a share, \texttt{na} values excluded.
\texttt{validate(rec, \dots)} checks a reading file (\cref{sec:formats}).

\subsection{\texttt{pineau\_rl.fmt}}

Numbers come out of this module the way the texts write them. \texttt{half\_up(x, nd)} rounds half away
from zero, as taught at school ($1.15 \to 1.2$), whereas Python's \texttt{round} rounds half to
even. The float goes through its shortest representation, so that 1.15 stays 1.15. \texttt{pct},
\texttt{num}, \texttt{ci} and \texttt{pts} produce a share as a percentage, an integer with a
thousands separator, an interval in whole points and the signed bounds of a difference, in French or
English style.

\section{File formats}
\label{sec:formats}

\paragraph{A simulation run} The file \file{sim/results/<algo>/<task>/seed<k>.csv} holds one
run, that is one program, task and random seed, with one line per evaluation (\cref{tab:format-essai}).
\begin{table}[htbp]
\centering
\small
\caption{The first two lines of the file of one simulation run: PPO PyTorch on \task{HalfCheetah}, random seed 0. Under each heading, the name of the column in the CSV file. The empty cell is a measurement not made.}
\label{tab:format-essai}
\begin{tabular}{@{}rrrrr@{}}
\toprule
Step & Mean return & Standard deviation & Mean return, modified physics & Elapsed seconds \\
{\footnotesize\texttt{step}} & {\footnotesize\texttt{eval\_mean}} & {\footnotesize\texttt{eval\_std}} & {\footnotesize\texttt{perturbed\_mean}} & {\footnotesize\texttt{wall}} \\
\midrule
\num{50000} & \num{329.4} & \num{26.6} & & \num{85.5} \\
\num{100000} & \num{690.4} & \num{71.0} & & \num{188.4} \\
\bottomrule
\end{tabular}
\end{table}
\texttt{eval\_mean} and \texttt{eval\_std} give the mean and the standard deviation of the returns of
the ten evaluation episodes, \texttt{perturbed\_mean} the mean return under modified physics (empty
when not measured), \texttt{wall} the elapsed seconds. A run is finished when its last step reaches
the program's budget. \Cref{tab:format-essai} gives the first two lines of PPO PyTorch on
\task{HalfCheetah}, random seed 0, rounded to one decimal.

\paragraph{A GRPO run} One JSONL file per model, reward and random seed
(\cref{chap:07-modeles-de-langue}). Evaluation lines contain \texttt{step}, \texttt{greedy},
\texttt{pass1}, \texttt{pass8} and \texttt{truncated}. Training lines, every ten steps, contain
\texttt{train\_right}, \texttt{truncated}, \texttt{loss}, \texttt{grad\_norm}, \texttt{retries} and
\texttt{wall}. These files keep no verdict per problem. The rerun of deviation~20 adds, for each run, a file
\file{seed<seed>.items.jsonl.gz} that keeps the verdict, the length and the answer for every problem at
every evaluation. Its log writes one line per step.

\paragraph{A paper reading} One JSON file per paper and per reader, stored under
\file{study/codes\_v2/} in the reader's directory. For each criterion, it holds the value, an exact
quotation from the paper and its page, or the statement \enquote{absent after reading the
experiments, appendices and links} with a null page (\cref{chap:A-grille}):
\begin{lstlisting}
{"id": "K001", "coder": "H", "title": "...",
 "criteria": {
  "Q1":  {"value": "yes", "quote": "We evaluate on six MuJoCo tasks", "page": 6},
  "X1":  {"value": 5, "quote": "averaged over 5 random seeds", "page": 7},
  "Y11": {"value": "yes", "kind": "se",
          "quote": "shaded areas denote one standard error", "page": 7}},
 "doubts": "..."}
\end{lstlisting}

\paragraph{The validator} \texttt{pineau\_rl.grid.validate} returns the list of problems in a reading
file, empty when the file is valid. It checks the schema: every criterion present, values
\texttt{yes}, \texttt{no} or \texttt{na}, integers or \texttt{null} for X1, X6 and Y3, an X5 family
among \texttt{llm}, \texttt{classic} and \texttt{theory}, a \texttt{kind} field when Y11 is yes, a
quotation for every criterion. It checks the scope rules of deviations 3 and 14 of the protocol:
\begin{itemize}
  \item Q1 set to no forces \texttt{na} everywhere;
  \item Y2 and Y11 are \texttt{na} if and only if X2 is no;
  \item Y4 is \texttt{na} when Y3 is 1;
  \item Y5 is \texttt{na} if and only if Q4 is \texttt{na}.
\end{itemize}
Given a function that searches the text of the PDF, the validator finally checks that every
quotation appears there. The function
\texttt{segments} first removes the reader's notes in square brackets, splits a quotation stitched
together with \enquote{[...]} into pieces and ignores pieces shorter than twelve letters. The script
\file{study/check\_codes.py} applies this check with the PDFs. The reader's sheet asks each reader to
run it on their files until it reports zero problems.

\section{Every number checked against its file}
\label{sec:verification}

Two scripts check the published texts. They read no number typed by hand: they recompute it.

\paragraph{\file{docs/check\_numbers.py}} For each numerical claim of the summary document, the
script recomputes the value from a result file, formats it as each language writes it
(\texttt{pineau\_rl.fmt}), then searches the text for the resulting string, with spaces and line
breaks normalised. A few survey counts that the analysis scripts do not write are recounted from
the readings, with the same rules. Assertions protect the qualitative facts, for example the name of
the task on which the leader changes under modified physics. The script checks 188 claims in the two
languages and stops at the first missing number: \file{docs/build.sh} therefore does not compile a
document in which a figure has drifted.

\begin{lstlisting}[language=Python,caption={One claim of \file{docs/check\_numbers.py}: the expected string is built from the data, never typed.},captionpos=b]
add("agreement",
    f"{same} décisions sur {total} ({pct(same / total)})",
    f"{same} decisions out of {total} ({pct(same / total, lang='en')})")
\end{lstlisting}

\paragraph{\file{article/check\_article.py}} The script of the blog article adds a coverage check.
Each claim builds, from a result file, the sentence that the two languages write, then searches the
text for it. The few numbers without a file (a slide, a cited paper, a sample size) are searched for
in the summary document. The script then removes from the text all checked sentences and the proper
names that contain digits: any number left over is a failure. The bibliography lines are confronted
with the file \file{docs/refs.bib}. The article has 161 checked claims in each language.

\section{The targets of the \file{Makefile}}
\label{sec:make}

The \file{Makefile} trains nothing. Training goes through \file{sim/run.py}, \file{sim/drive.sh} and
\file{sim/llm/grpo.py}.

\begin{description}[style=nextline,leftmargin=1.5em]
  \item[\texttt{make install}] installs the package in editable mode in the analysis environment,
        with no dependency downloaded: they come from \file{envs/analysis.lock}.
  \item[\texttt{make analysis}] reruns every analysis script, in dependency order, through three
        targets that also run alone. \texttt{analysis-study} runs the survey scripts, among them
        those of deviations 30 to 36. \texttt{analysis-sim} runs the seventeen simulation scripts at ten
        random seeds, among them the strict count of E2 and the test of the \task{Hopper} reversal, then
        the statistical analysis of the official Zoo pipeline, the variants of deviations 12 and 20
        (twenty random seeds, replication of E4 on random seeds 10 to 19, equal budget) and the eight
        analyses of deviations 22 to 31, which read the files of the earlier ones. \texttt{analysis-llm}
        runs \file{summary.py} at three and at five random seeds, then \file{logs\_stats.py}
        (deviation~19). The variants are selected through environment variables (\texttt{SEEDS},
        \texttt{SUFFIX}, \texttt{BUDGET}). Two analyses of deviation~33, \file{e6\_enjoy.py} and
        \file{e6\_hub.py}, need the Zoo environment and the policies: no target reruns them, as with
        \file{e6\_official\_eval.json}.
  \item[\texttt{make hyperparams}] reads the settings of the five programs back from the libraries of the
        training environment and writes \file{sim/analysis/hyperparams.json} (deviation~28). This target stays
        outside \texttt{verify}, which runs in the analysis environment.
  \item[\texttt{make analysis-suite}] runs the three analyses written in advance for the experiments
        not yet finished (deviations 32 to 34). Their JSON files enter the verification only once their
        training runs have finished.
  \item[\texttt{make check}] runs the two checking scripts of \cref{sec:verification}.
  \item[\texttt{make test}] runs the unit tests with pytest.
  \item[\texttt{make lint}] runs ruff and mypy on the package and the tests.
  \item[\texttt{make verify}] chains \texttt{analysis}, \texttt{check} and \texttt{test}, then asks
        git whether every JSON result file has come back byte for byte.
  \item[\texttt{make verify-head}] does the same in a fresh export of the last commit, without the
        policies, the Zoo runs or the training logs, which no analysis reads. This target leaves
        untouched a working copy in which runs are still writing.
  \item[\texttt{make figures}] regenerates the figures of the summary document
        (\file{docs/figures/build.py}).
  \item[\texttt{make docs}] runs \texttt{check}, then builds the summary document
        (\file{docs/build.sh}).
  \item[\texttt{make report}] builds this report (\file{report/build.sh}): tables generated from
        the JSON files, figures, then LuaLaTeX and biber, in French and in English.
\end{description}

\begin{lstlisting}[language=bash,caption={Recomputing every number from the committed results, without training.},captionpos=b]
# Create the analysis environment
python3.11 -m venv sim/env-an
# Install the packages of the analysis and of the tests
sim/env-an/bin/pip install -r envs/analysis.txt -r envs/test.txt
# Analyses, numbers of the texts, tests, byte-for-byte diff
make verify
# The same in a fresh export of the last commit
make verify-head
\end{lstlisting}

\section{Continuous integration}
\label{sec:ci}

The file \file{.github/workflows/verify.yml} replays the verification on every push and every pull
request, on two platforms. The GitHub actions it uses are pinned to a commit hash, since a tag can be moved. Each
job installs the analysis environment from \file{envs/analysis.lock}, with every package checked
against its SHA-256 hash (\texttt{pip install -{}-require-hashes}), then installs the package, runs
the tests, reruns all the analyses, checks the numbers of the texts and compares every result file
with the commit using \file{tools/compare\_results.py}.
\begin{itemize}
  \item On Ubuntu 24.04 (x86-64), with Python 3.11.18, the version we use, the comparison
        accepts a relative difference of $10^{-12}$ on floats. It requires equality of counts,
        verdicts, rankings and all strings. The $p$-values move in their fifteenth or sixteenth digit
        with the mathematics library.
  \item On macOS 15 (arm64), a system of the machine family that produced the results (a Mac under
        macOS 27.2, \cref{sec:log-calcul}), the comparison requires exact equality.
\end{itemize}
A third job builds the image of the \file{Dockerfile} and runs the tests and the number check in it.
The image starts from \texttt{python:3.11-slim}, pinned by hash, installs the same locked file and
mounts the repository at launch instead of copying it, so the same image serves any commit.

\section{Tests}
\label{sec:tests}

The \file{tests/} directory holds eight test files and one file of shared helpers
(\file{conftest.py}). The tests import the package and the thin layer \file{sim/analysis/runs.py},
never an analysis script, which computes as soon as it is imported. The tests combine three kinds of
checks: values published in the literature, an independent computation (statsmodels, SciPy,
NumPy) and reference values on synthetic data, frozen to a relative $10^{-12}$. A reference value
that changes announces a result file that changes.

\begin{description}[style=nextline,leftmargin=1.5em]
  \item[\file{test\_stats.py}] Wilson against a published value, against statsmodels and on its
        properties (the interval contains the share and stays within $[0,1]$, it is symmetric, it
        widens at 99\%); Newcombe against a published value and rebuilt from the two Wilson
        intervals; Fisher against SciPy; Holm computed by hand and its properties; Cochran-Armitage
        equal to $\sqrt N$ times the Pearson correlation between score and outcome; Cohen's kappa
        and its list of categories; the probability of improvement.
  \item[\file{test\_bootstrap.py}] the interquartile mean; the rliable intervals, identical under a
        fixed global random seed and dependent on that random seed alone; the bootstrap probability of
        improvement, which uses only its own generator.
  \item[\file{test\_henderson.py}] the three readings of Henderson's statistic against SciPy; the
        $p$-value of the mean $t$ with the mean Welch degrees of freedom; the computation of the
        script, which ignores the spread between random seeds: on six synthetic random seeds, it
        reaches a $t$ of 9 where a test over the random seeds sees nothing.
  \item[\file{test\_runs.py}] the format of a run: writing then reading, missing or unfinished runs
        returned as \texttt{None}, malformed files that raise an error, final scores and reference
        grid.
  \item[\file{test\_grid.py}] the two value semantics, the conversion of Y3, the harmonisation, the
        validator (valid file, each problem reported, scope rules, quotations) and the splitting of
        quotations.
  \item[\file{test\_formatting.py}] rounding half away from zero, percentages, thousands, intervals
        and the minus sign of both languages.
  \item[\file{test\_pipelines.py}] the reading of runs by \file{sim/analysis/runs.py} and the GRPO
        summary from end to end, on three synthetic random seeds: an unfinished run is set aside, while
        the mean gain and the bootstrap interval match their expected values.
  \item[\file{test\_draw.py}] the 2026 draw comes back from its rule (\file{study/draw\_2026.py}) and
        depends on its random seed.
\end{description}

\section{Frozen environments}
\label{sec:envs}

Four Python environments ran the project, all under Python 3.11.18. Each file in \file{envs/} is the
\texttt{pip freeze} of one of them, with the Python version and the
platform at the top. The \file{*.lock} files add the SHA-256 hash of each package.
\Cref{chap:D-environnements} gives the versions. The training environment (\file{envs/train.txt})
serves \file{run.py} and \file{dose.py}, the analysis environment (\file{analysis.txt}) the analyses,
the verification and continuous integration, the Zoo environment (\file{zoo.txt}) the official
pipeline of experiment E6, the language-model environment (\file{llm.txt}) GRPO and the figures
of the summary document. A fifth file, \file{test.txt}, adds pytest to the analysis environment.

\section{The trained policies}
\label{sec:release}

The final policies let anyone re-evaluate an agent without retraining it, as the graded
re-evaluation in modified physics did (\file{sim/dose.py}). Those of PPO, TRPO and experiment E6 are
in the repository. The weights of SAC and CrossQ, too heavy for git, are attached to the GitHub
release \texttt{policies-v1}, as tar archives listed with their SHA-256 hash in a \file{SHA256SUMS}
file. The script \file{sim/fetch\_models.py} downloads each archive, checks its hash, unpacks it
under \file{models/}, then checks each policy file against \file{sim/models/MANIFEST.sha256}, the
hashes taken on the machine that trained them. The downloads run one after
another, separated by a random pause of 1.5 to 4 seconds. A refused request is retried after a
wait drawn at random between zero and a ceiling that doubles with each failure (exponential backoff
with full jitter), five attempts at most.

\begin{lstlisting}[language=bash,caption={Fetching and checking the policies of SAC and CrossQ.},captionpos=b]
# Download, unpack, check everything
cd sim && python3 fetch_models.py
# Only check the policies already present
cd sim && python3 fetch_models.py --check
\end{lstlisting}

\section{Hardware and compute}
\label{sec:log-calcul}

All runs ran on one laptop, a Mac with an Apple M2 Max chip under macOS 27.2 (arm64), shared
with other projects. Each simulation training runs on a single thread. The number of jobs in
parallel changed over the project: four cores at the start, up to eight granted every hour according
to the load of the machine (deviation 15 of the simulations), then two cores for the whole project
(deviation 19). This sharing leaves the results untouched as long as a run
depends on its random seed alone: the relaunches check it for PPO, while deviation 32 is still
testing it for SAC and CrossQ (\cref{sec:registres}). \Cref{sec:calcul} describes the launchers, the watchdog and the job queue.
\Cref{tab:calcul} gives, per program and per task, the number of runs, the number of steps, the
median wall-clock time of a run and the sum of these times.

The GRPO runs have very different durations depending on the device. According to the \texttt{wall}
field of their logs at step 100, a Qwen run took 42 to 52 minutes on GPU and 4.3 to 11.4 hours on
CPU, a Llama run 69 to 148 minutes on GPU. These spreads also reflect the varying load of the
machine. From these logs, the rerun of deviation~20 plans about 58 GPU hours for its forty runs, to be
run on an idle machine.

\section{What is not in the repository}
\label{sec:absent}

The PDFs of the papers read are protected by copyright. The script \file{study/fetch.py} downloads
them again from the public sources (\texttt{cd study \&\& python3 fetch.py final\_2026.json pdf}).
Pineau's slides are not in the repository for the same reason: they are on the websites of ICLR 2018
and NeurIPS 2018. The file \file{slides-en.md} describes them one by one. The logs of the working
sessions with the agents that read the papers are not public. The exact instructions, the model and
the dates of each reader, reconstructed from these logs, are public (\file{study/readers/}).

\section{Licence and citation}
\label{sec:licence}

The code, the protocols, the derived data and the documents are published under the BSD 3-Clause
licence, the licence of scikit-learn (file \file{LICENSE}, Copyright (c) 2026, Warith Harchaoui). The
file \file{CITATION.cff} gives the reference, which GitHub offers in APA and BibTeX.
